
% ============================================================================
\clearpage
\section*{v4.10.0-r64 Source-Only: Record of the Tests H105 and H128--H136, and Precisions from the Audit of the Physics Article}
\addcontentsline{toc}{section}{v4.10.0-r64 Source-Only: Record of the Tests H105 and H128--H136, and Precisions from the Audit of the Physics Article}
\label{sec:v4100-r64-record-audit}
% ============================================================================

\paragraph{Purpose.}
On 26 November 2026 the stand-alone physics article ``One Object, Three Frames'' (version 1.0.0, doi:10.5281/zenodo.24539644) was audited
sentence by sentence against this Companion, the test register and the research notes written after r60.  Where the corrected article now
refers to a result, that result has to stand in the Companion.  The layers r62 and r63 collect H138--H156.  This layer collects the remaining
tests that the article cites, H128--H136 of 25 November 2026 and the open prediction H105 of 16 November 2026, each with the status recorded in
its result note.  It also records the precisions to earlier layers that the audit found.  It adds no new computation.  No test status and no
guard changes.  Earlier layers are not rewritten.  The test H137 (tape circuit on superpositions) belongs to the QAM line and is not collected here.

\noindent\textbf{Source-only status.}
This v4.10.0-r64 layer is source-only.  It does not build a PDF, does not assign a new DOI,
and does not authorize a public release.  The v5.4.0 Version DOI
\href{https://doi.org/10.5281/zenodo.24500009}{10.5281/zenodo.24500009} and the Concept DOI
\href{https://doi.org/10.5281/zenodo.19210728}{10.5281/zenodo.19210728} are unchanged by it.

\subsection*{v4.10.0-r64 results recorded}
\label{subsec:v4100-r64-results}

\paragraph{Theoremlet v4.10.0-r64.1 (what the step to the sedenions leaves of the family symmetry; H128, k283; partly not blind).}
For the Hermitian \(3\times3\) matrices over the Cayley--Dickson algebras, \(\dim\operatorname{Der}H_3(A_n)=3,\,8,\,21,\,52\) for
\(n=0,\dots,3\) (control) and \(\dim\operatorname{Der}H_3(\mathbb S)=22\), with the structure \(g_2\oplus su(3)\) (numerical rank with a gap of
\(1.5\cdot10^{13}\); the exact rank was not obtained).  The pre-registered value \(17=14+3\) is \textbf{refuted} (H128a).  The \(su(3)\) is the
commutant of \(g_2\); it contains the family \(so(3)\).  For a charge direction \(u\) the ratio of the quadratic mass forms of the spin-1 and the
spin-2 part is \(r(u)=1-(1-u_8^2)/10\) exactly (Wolfram Language): \(r=9/10\) for \(u=e_7\), H128b \textbf{confirmed}; \(r=1\) for the doubling
unit \(u=e_8\), H128b \textbf{refuted} for that direction.  The sign with which this form enters a potential is open and is not set.  The step
\(SO(3)\to A_4\) does not come from the tower (as disclosed in advance; it is the charge per slot of r60).

\paragraph{Theoremlet v4.10.0-r64.2 (octonionic colour-stable carriers exist only perpendicular to the doubling unit; H130, k258; partly not blind).}
Demand a charge direction \(u\in\operatorname{span}(e_7,e_8,e_{15})\) and a closed colour-stable subalgebra of \(\mathbb S\) on which
\(su(3)_c\) acts as \(1\oplus1\oplus3\oplus\bar3\).  All 600 random starts end in one of three classes: \(u=e_7\) and
\(u=-\tfrac12e_7\pm\tfrac{\sqrt3}{2}e_{15}\) (alternative and norm multiplicative, the three octonion copies of r62.1; \(u_8=0\) in 383 of 383
solutions), and \(u=e_8\), where the closed algebra is \emph{not} alternative and \emph{not} norm multiplicative (twelve non-zero alternator
components, Wolfram Language exact).  H130a and H130b \textbf{confirmed}.  Under this demand the ratio \(9/10\) of r64.1 is the one that applies.
\emph{Consequence for the register:} the placeholder CT-003 (charge axis \(e_8\), r13.1) is not supported by this finding: at \(u=e_8\)
the closed colour-stable algebras found are not alternative.  Route A of r13 (two generations) and the relative CP sign of r59.6 keep their status ``conditional on
CT-003''; whether they survive on \(u=e_7\) is not computed.  CT-003 stays a placeholder.

\paragraph{Theoremlet v4.10.0-r64.3 (mixing from the copies: one class, and the frame rule for comparisons with data; H131, H134, k260--k263; partly not blind).}
With \(\mathbb Z_3\) for charged leptons, \(\mathbb Z_2\) with a generalised CP from the charge for Majorana neutrinos and ``copy \(k\) on slot
\(k\)'' (all settings), the 48 admissible combinations give one class: a trimaximal column with \(J=0\), \(\delta\in\{0,\pi\}\) and
\(\sin^2\theta_{23}=0.395\) or \(0.605\) at \(\sin^2\theta_{13}=0.022\) (found after the fact).  H131a \textbf{confirmed}, H131b (CP violation with
\(\mu\tau\) reflection) \textbf{refuted}.  The result note sets this column beside measured values without a declared map to the matter frame.
The author's objection of 25 November 2026 is recorded as a rule: \emph{before a statement ``disfavoured by data'', the prediction is to be
transferred to the matter frame.}  Frame maps that commute with the group (\(z\mapsto iz\), \(z\mapsto\bar z\)) cannot move a column fixed by a
residual symmetry; the only map in the Companion that distinguishes the slots is the mask \(D=\exp(\varepsilon W_l)\) with
\(\varepsilon=\operatorname{artanh}\beta_{\rm sym}=0.6174\) (H45).  With it the TM1 class has 56 hits compatible with the JUNO value of
\(\theta_{12}\), only if the mask also acts on the neutrinos (H134a, H134c \textbf{confirmed}); \(\theta_{23}\) and \(\delta\) are not sharp
(widths \(0.635\) and \(0.932\); H134b \textbf{refuted}).  This is compatibility, not a prediction.  The verdict ``stands badly against the data''
for the H131 column therefore holds only for the untransferred column and is not a verdict on the construction.

\paragraph{Theoremlet v4.10.0-r64.4 (a recursive field equation on the tower fixes no scale; H129, H132, H133, H135, H136, k290--k305; partly not blind).}
An ansatz \(L_{n+1}[\psi+\chi\ell]=L_n[\psi]+\Delta_n[\psi,\chi]\) with one new term per step (area, commutator, associator, norm defect) has the
symmetry of each step; its signs are free parameters (exploration k290).  It lies outside the setting W1/W2 of v4.9.0-r81--r86 and outside the
list \(\mathcal K\) of r56.  Tested: (i) H129: one-loop flow generates a scale on \(67.8\,\%\) of the grid (carries partly), but not in the
appended half field (\(29.5\,\%\), equal by exchange symmetry to the share in the pre-stage); \textbf{not confirmed}.  (ii) H132: on the HC
direction no gauge field stays massless (\(\det M^2=4g_L^2g_R^2AB\)); \textbf{refuted}; for \(g_L=g_R\) the ratio is
\(m_-/m_+=\beta_{\rm sym}^2\) with \(\beta_{\rm sym}\) inserted.  (iii) H133: the tower fixes the form of the coupling of the \(HS(3)\) appendix
(one mixed term \(\kappa AC\), four degenerate fields) but not its value; \textbf{not confirmed}.  (iv) H135, H136: \(\kappa\approx1.2\) is neither
a fixed point of the flow, nor an extremum of the vacuum depth, nor held by the back-reaction through the HC angle; both \textbf{refuted}.  The
numerical hit \(\kappa=(\gamma_L+1/\gamma_L)/2=1.2004\) (k303) uses \(\gamma_L\) as input, lies at chance level among the candidates searched and
is recorded as no evidence.  Consequence: as in v4.9.0-r84 and r56, no scale and no value
follows; OP1 and the scale stay open.

\paragraph{Remark v4.10.0-r64.5 (H105 is an assigned, open prediction).}
The statements ``H105 stays unassigned'' (r18) and ``no test'' (r24.4) are superseded.  On 16 November 2026 the author chose the setting V1
(every crossing of the horizon is one T quantum) so that the crossing reading of r24.4 has a measurable consequence.  H105, pre-registered
(note \texttt{Vorab\_\allowbreak Notiz\_\allowbreak H105\_\allowbreak Flaechenschritte\_\allowbreak Echos.md}, script k227): area steps
\(\Delta A=4\ln2\,\ell_P^2\), absorption lines at \(\omega_n=n\,\tfrac{\ln2}{2\pi}\kappa+m\Omega_H\), hence echoes after the ringdown; to be
decided by gravitational-wave analyses appearing after 16 November 2026, deadline 31 December 2028 (confirmed at Bayes factor \(\ge10\) or
\(\ge3\sigma\), refuted at exclusion for one event).  Status: open.  The value \(\alpha=4\ln2\) is not specific to the construction (any counting
with \(j=\tfrac12\) only gives it), and whether such a quantisation produces observable echoes is disputed in the literature.

\subsection*{v4.10.0-r64 precisions to earlier layers}
\label{subsec:v4100-r64-precisions}

\paragraph{Remark v4.10.0-r64.6 (numbers).}
\begin{enumerate}[leftmargin=2em]
\item[(a)] r08.3(a): with the SU(2) counting (\(\gamma_L=2.151501\)) the \(m_W\) relation gives \(82.69\) GeV (stored output of k218:
\(82.69234\)), not \(82.68\) GeV; the latter came from the rounded \(\gamma_L\approx2.15\).  \(\Omega_M=0.365\) is unchanged.
\item[(b)] Two values of \(\gamma_L\) appear in the Companion: \(1.864700\) from \(\gamma_M=0.23753296\) and \(1.86444\) from the rounded input
\(0.2375\) of the main body (v4.9.0-r68).  The entries \(1.86444\) (v4.9.0-r24.3) and \(1.8645\) (v4.9.0-r84.2(b)) are the rounded-input value;
with the unrounded input both read \(1.8647\).  No statement depends on the difference.
\item[(c)] Koide: with \(m_\tau=1776.93\pm0.09\) MeV, \(K=0.6666645\pm0.0000051\), that is \(0.43\sigma\) below \(\tfrac23\) (r63.5).  The
``within \(0.005\,\%\)'' of the main text is correct but loose by a factor of about fifteen.  H102 (\(m_\tau=1777.134\) MeV) uses
\(\rho^2=0.500014\) taken from \(m_e/m_\mu\), not the overlap \(\tfrac12\); with exactly \(\tfrac12\) the value is \(1776.969\) MeV.  The two
statements do not contradict each other: the deviation that H102 explains is not significant at present.
\end{enumerate}

\paragraph{Remark v4.10.0-r64.7 (which map exchanges the frames).}
\begin{enumerate}[leftmargin=2em]
\item[(a)] Remark r29.4 writes ``the frame T is the \(J_B\) image with \(s\mapsto-\bar s\)''.  By the correction v4.9.0-r28 the frame T of
Convention F is the image under the duality rotation \(D\colon z\mapsto iz\), \(s\mapsto-s\).  Both maps reverse \(\operatorname{Re}s\), and the
argument of r29.4 (a boost preserves \(s\), so \(\gamma_L\) is the factor of \((M,\bar M)\), not of \((M,T)\)) holds for either.  r29.4 is to be
read with \(D\).
\item[(b)] The exchange of r63.2 and r63.7, \(J\colon\psi\mapsto U\bar\psi\), is antilinear.  On \(\operatorname{span}(1,U)\cong\mathbb C\) it is
\(z\mapsto i\bar z\), the map \(J_B=D\circ\vartheta\), not \(D\).  \(D\) itself acts as \(\theta\mapsto\theta+\pi/2\) (H149b) and has no fixed
angle; \(J\) acts as \(\theta\mapsto\pi/2-\theta\) and fixes \(\theta=\pi/4\).  \(J\) exchanges matter-type and tachyon-type states
(\(\operatorname{Re}\Psi^2\mapsto-\operatorname{Re}\Psi^2\)); calling it ``the exchange of matter frame and tachyon frame'' is the declared
reading of k297 (\(\mu\mapsto1/\mu\)), not the tachyon map of Convention F.  \(J\) is not time reversal \(\vartheta\colon z\mapsto\bar z\) either.
\item[(c)] At the fixed point \(\operatorname{Re}\Psi^2=0\) but \(\Psi^2=U\ne0\).  ``Lightlike'' in r63.2 means \(\operatorname{Re}s=0\), the
boundary between the two classes; it is not the null class \(s=0\) and not the photon eigenframe \(P\) (\(b=0\)).
\end{enumerate}

\paragraph{Remark v4.10.0-r64.8 (scope of earlier statements).}
\begin{enumerate}[leftmargin=2em]
\item[(a)] r06.1, r07.2: Dorau and Much give arguments for the chain, for a real scalar field in coherent states at leading order.  That the
chain holds for the Maxwell field in frame M is an assumption of r07.2; r07.2 shows that frame T then gives the same equation.
\item[(b)] Main text, \texttt{prop:compensator}: the status note of v4.9.0-r03.4 corrects only the orthogonality clause.  The clause ``produces
the photon \(C(f)\in K_3\)'' is superseded by r23.2 and r11.2(f): the compensator image lies on the compensator axis R (\(\phi=1\)) and never
reaches \(P\).
\item[(c)] r30.6, r55.1: ``the numerical values in OP2--OP5 all go back to \(\beta_{\rm sym}\)'' holds for the inserted values.  The open
predictions H42, H100, H112 and H125 do not depend on \(\beta_{\rm sym}\) (r08.3, r12.5, r59), and the Koide overlap \(\tfrac12\) and the number
of generations are separate inputs.
\item[(d)] r19.4: the caveat ``to be confirmed in the PDFs before citation'' is lifted for the formulas quoted there; they were confirmed in the
full texts (Brannen, arXiv:1006.3114, eqs.~17, 56, 67, 69, 72; Foot, arXiv:hep-ph/9402242; \.Zenczykowski, arXiv:1210.4125).  Brannen attributes
the charged-lepton formula with \(2/9\) to G.~Rosen (Mod.\ Phys.\ Lett.\ A 22 (2007) 283).
\item[(e)] Numbering: in r25, r26 the names \(HS(6)\), \(HS(7)\), \(HS(8)\) denote Dixon copies built on \(\mathbb C\otimes\mathbb H\otimes\mathbb O\);
in r62, r63 and in this layer \(HS(n)\) denotes the Cayley--Dickson level \(A_n\) (\(HS(4)=\mathbb S\)).  The two are different towers.
\end{enumerate}

\subsection*{v4.10.0-r64 frames}
\label{subsec:v4100-r64-frames}

\begin{center}
\small
\begin{tabular}{@{}>{\raggedright\arraybackslash}p{0.34\textwidth}>{\raggedright\arraybackslash}p{0.20\textwidth}>{\raggedright\arraybackslash}p{0.38\textwidth}@{}}
\toprule
item & frame & comparison\\
\midrule
r64.1, r64.2 (derivations, carriers) & frame independent algebra & none with data\\
r64.3 (mixing) & algebraic column; M after the mask \(\varepsilon=\operatorname{artanh}\beta_{\rm sym}\) & with data only after the declared map (H45, H134)\\
r64.4 (field equation on the tower) & M; exchange \(\mu\mapsto1/\mu\) as declared reading & none with data\\
r64.5 (H105) & T counts crossings (setting V1) & gravitational-wave analyses, open\\
r64.7 & M, T & \(D\), \(J_B\) and \(\vartheta\) kept apart (v4.9.0-r28, r51, r81)\\
P & not used & ---\\
\bottomrule
\end{tabular}
\end{center}

\subsection*{v4.10.0-r64 register entries}
\label{subsec:v4100-r64-register}

\begin{center}
\small
\begin{tabular}{@{}>{\raggedright\arraybackslash}p{0.40\textwidth}>{\raggedright\arraybackslash}p{0.52\textwidth}@{}}
\toprule
statement & status after v4.10.0-r64\\
\midrule
H128a; H128b for \(e_8\); H131b; H132; H134b; H135; H136 & refuted\\
H128b for \(e_7\); H130a, b; H131a; H134a, c & confirmed, partly not blind\\
H129, H133 & not confirmed (carry partly)\\
H105 & open, deadline 31 December 2028\\
CT-003 (charge axis \(e_8\)) & placeholder; not supported by H130; dependent statements stay conditional\\
frame rule for comparisons with data & recorded (author, 25 November 2026)\\
scale and \(\beta_{\rm sym}\) from a field equation on the tower & not obtained\\
OP1--OP5 & open\\
\bottomrule
\end{tabular}
\end{center}

\subsection*{v4.10.0-r64 notes}
\label{subsec:v4100-r64-notes}

{\raggedright
Result and advance notes of H128--H136: \texttt{notes/\allowbreak research/\allowbreak predictions/} (H128, H130, H131, H134 and all advance
notes) and \texttt{notes/\allowbreak research/\allowbreak hilbert\_\allowbreak shift\_\allowbreak felder\_\allowbreak 2026\_\allowbreak 11\_\allowbreak 24/\allowbreak feldgleichung/}
(H129, H132, H133, H135, H136, explorations k290--k305), with the scripts, Wolfram cross-checks and actual outputs in the subfolders
\texttt{kollaps\_\allowbreak h128}, \texttt{ladungsrichtung\_\allowbreak h130}, \texttt{mischung\_\allowbreak h131}, \texttt{frame\_\allowbreak h134} and
\texttt{feldgleichung}.  H105: \texttt{notes/\allowbreak research/\allowbreak predictions/\allowbreak Vorab\_\allowbreak Notiz\_\allowbreak H105\_\allowbreak Flaechenschritte\_\allowbreak Echos.md}.
Audit of the article and list of its changes: \texttt{notes/\allowbreak research/\allowbreak v4\_\allowbreak 10\_\allowbreak 0\_\allowbreak r64\_\allowbreak nachtrag\_\allowbreak audit/\allowbreak README.md}.
\par}

\subsection*{v4.10.0-r64 ledger and guard update}
\label{subsec:v4100-r64-ledger-guard}

\[
R_{\rm chk}^{v4100\text{-}r64}=(1,1,0)
\]
for (status recorded; computations checked in the cited scripts; no new computation in this layer).

The layer spends no guard:
\[
R_{\rm glob}^{(v4100\text{-}r64,+)}=(0,1,1,1,1),
\qquad
\Gamma_{\rm guard}^{v4100\text{-}r64}=(1,1,1,1,1),
\qquad
\beta_4=1/256\ \text{unspent}.
\]
The layer does not imply exact finite \(C5\) zero closure, exact all-readout terminality, \(C6/T2\), a global photon-clock/RH theorem,
or global all-readout closure.  OP1--OP5, C5 and C6/T2 remain open.
